%% ma: L10-B13-0037
%% lop: 10
%% chuong: 5
%% bai: 13
%% dang: 10.13.2
%% loai: TN4
%% muc: TH
%% dap_an: "C"
%% nguon: "chuyen-de-cac-so-dac-trung-cua-mau-so-lieu-khong-ghep-nhom-toan-10-chuong-trinh-moi.pdf (GV. Phan Nhật Linh) (P02), Bài 02 Dạng 1 Phần I Câu 28"
%% kien_thuc: "Lớp 10 Bài 13 (số trung bình)"
%% trang_thai: chinh-thuc
%% ghi_chu: "Đáp án gốc đánh dấu A (15,20) sai, lời giải gốc ra 15,23 (C). Tự duyệt theo yêu cầu giáo viên 22/10/2026."
\begin{ex}
Có $100$ học sinh tham dự kì thi học sinh giỏi môn Hóa học (thang điểm 20). Kết quả như sau:
\begin{center}
\begin{tabular}{|c|c|c|c|c|c|c|c|c|c|c|c|}\hline
Điểm & $9$ & $10$ & $11$ & $12$ & $13$ & $14$ & $15$ & $16$ & $17$ & $18$ & $19$\\\hline
Tần số & $1$ & $1$ & $3$ & $5$ & $8$ & $13$ & $19$ & $24$ & $14$ & $10$ & $2$\\\hline
\end{tabular}
\end{center}
Số trung bình của mẫu số liệu là
\choice{$\overline x=15{,}20$}{$\overline x=15{,}21$}{\True $\overline x=15{,}23$}{$\overline x=15{,}25$}
\loigiai{Tổng điểm: $9\cdot1+10\cdot1+11\cdot3+\cdots+18\cdot10+19\cdot2=1523$.

$\overline x=\dfrac{1523}{100}=15{,}23$. Chọn C.}
\end{ex}
